% Options for packages loaded elsewhere
\PassOptionsToPackage{unicode}{hyperref}
\PassOptionsToPackage{hyphens}{url}
%
\documentclass[
]{article}
\usepackage{amsmath,amssymb}
\usepackage{iftex}
\ifPDFTeX
  \usepackage[T1]{fontenc}
  \usepackage[utf8]{inputenc}
  \usepackage{textcomp} % provide euro and other symbols
\else % if luatex or xetex
  \usepackage{unicode-math} % this also loads fontspec
  \defaultfontfeatures{Scale=MatchLowercase}
  \defaultfontfeatures[\rmfamily]{Ligatures=TeX,Scale=1}
\fi
\usepackage{lmodern}
\ifPDFTeX\else
  % xetex/luatex font selection
\fi
% Use upquote if available, for straight quotes in verbatim environments
\IfFileExists{upquote.sty}{\usepackage{upquote}}{}
\IfFileExists{microtype.sty}{% use microtype if available
  \usepackage[]{microtype}
  \UseMicrotypeSet[protrusion]{basicmath} % disable protrusion for tt fonts
}{}
\makeatletter
\@ifundefined{KOMAClassName}{% if non-KOMA class
  \IfFileExists{parskip.sty}{%
    \usepackage{parskip}
  }{% else
    \setlength{\parindent}{0pt}
    \setlength{\parskip}{6pt plus 2pt minus 1pt}}
}{% if KOMA class
  \KOMAoptions{parskip=half}}
\makeatother
\usepackage{xcolor}
\usepackage{longtable,booktabs,array}
\usepackage{calc} % for calculating minipage widths
% Correct order of tables after \paragraph or \subparagraph
\usepackage{etoolbox}
\makeatletter
\patchcmd\longtable{\par}{\if@noskipsec\mbox{}\fi\par}{}{}
\makeatother
% Allow footnotes in longtable head/foot
\IfFileExists{footnotehyper.sty}{\usepackage{footnotehyper}}{\usepackage{footnote}}
\makesavenoteenv{longtable}
\usepackage{graphicx}
\makeatletter
\newsavebox\pandoc@box
\newcommand*\pandocbounded[1]{% scales image to fit in text height/width
  \sbox\pandoc@box{#1}%
  \Gscale@div\@tempa{\textheight}{\dimexpr\ht\pandoc@box+\dp\pandoc@box\relax}%
  \Gscale@div\@tempb{\linewidth}{\wd\pandoc@box}%
  \ifdim\@tempb\p@<\@tempa\p@\let\@tempa\@tempb\fi% select the smaller of both
  \ifdim\@tempa\p@<\p@\scalebox{\@tempa}{\usebox\pandoc@box}%
  \else\usebox{\pandoc@box}%
  \fi%
}
% Set default figure placement to htbp
\def\fps@figure{htbp}
\makeatother
\setlength{\emergencystretch}{3em} % prevent overfull lines
\providecommand{\tightlist}{%
  \setlength{\itemsep}{0pt}\setlength{\parskip}{0pt}}
\setcounter{secnumdepth}{-\maxdimen} % remove section numbering
\usepackage{bookmark}
\IfFileExists{xurl.sty}{\usepackage{xurl}}{} % add URL line breaks if available
\urlstyle{same}
\hypersetup{
  pdftitle={03. IoT Architecture},
  hidelinks,
  pdfcreator={LaTeX via pandoc}}

\title{03. IoT Architecture}
\author{}
\date{}

\begin{document}
\maketitle

\section{IoT System Architecture}\label{iot-system-architecture}

\subsection{Communication Paradigms}\label{communication-paradigms}

\begin{itemize}
\tightlist
\item
  \textbf{Client/Server}: sources explicitly know the destination of
  their messages.
\item
  \textbf{Publish/Subscribe}: sources publish to a logical channel;
  interested destinations subscribe (one-to-many, no explicit sessions).
\item
  Typical IoT traffic flows:

  \begin{itemize}
  \tightlist
  \item
    \textbf{Asynchronous events} (uplink)
  \item
    \textbf{Periodic or on-demand measurements} (uplink)
  \item
    \textbf{Commands} from processing platforms (downlink)
  \end{itemize}
\end{itemize}

\begin{center}\rule{0.5\linewidth}{0.5pt}\end{center}

\subsection{Properties of an IoT
System}\label{properties-of-an-iot-system}

\begin{itemize}
\tightlist
\item
  \textbf{Trustworthiness}: availability, resilience, confidentiality,
  integrity, safety.
\item
  \textbf{Architectural}: composability, modularity, heterogeneity,
  scalability, legacy support, shareability, unique identification.
\item
  \textbf{Functional}: accuracy, auto-configuration, context-awareness,
  big data management, discoverability, real-time, self-description,
  compliance.
\end{itemize}

\begin{center}\rule{0.5\linewidth}{0.5pt}\end{center}

\subsection{Standard Reference Model}\label{standard-reference-model}

\begin{itemize}
\tightlist
\item
  \textbf{Physical Entity Domain (PED)}: monitored physical objects.
\item
  \textbf{Sensing \& Controlling Domain (SCD)}: sensors, actuators
  bridging physical ↔ cyber world.
\item
  \textbf{Operations \& Management Domain (OMD)}: OSS + BSS for
  provisioning, assurance, optimization.
\item
  \textbf{Resource Access \& Interchange Domain (RAID)}: service
  interfaces with access policies.
\item
  \textbf{Application \& Service Domain (ASD)}: end-user services via
  APIs or portals.
\item
  \textbf{User Domain (UD)}: human/digital users via devices (PCs,
  smartphones, control panels, smart glasses)
\end{itemize}

\pandocbounded{\includegraphics[keepaspectratio]{/home/alessandrovulcu/Documenti/Università/ICT/Pasted image 20251001094808.png}}

\begin{center}\rule{0.5\linewidth}{0.5pt}\end{center}

\subsection{Network Architectures}\label{network-architectures}

\begin{itemize}
\tightlist
\item
  \textbf{One-level}: endpoints use full IP stack (e.g., cellular).
\item
  \textbf{Two-level (1b)}: endpoints resource-constrained, gateway
  manages protocols (LPWA, Low Power Wide Area: LoRa, SigFox, NB-IoT).
  \emph{Target of the course}
\item
  \textbf{Two-level (1c)}: endpoints support TCP/IP, traffic managed by
  router (Ethernet/Wi-Fi)
\end{itemize}

\begin{center}\rule{0.5\linewidth}{0.5pt}\end{center}

\subsection{WANs}\label{wans}

\begin{itemize}
\tightlist
\item
  TCP/IP networks, shared among users.
\item
  Public by nature.
\item
  \textbf{VPNs}: reserved bandwidth/security for defined groups, with
  private IP addressing
\item
  Users outside of a group cannot communicate directly with users of the
  group.
\item
  Use private IP addressing and private routing protocol instances.
\end{itemize}

\begin{center}\rule{0.5\linewidth}{0.5pt}\end{center}

\subsection{Topologies}\label{topologies}

\begin{itemize}
\tightlist
\item
  \textbf{Bus}: every nodes taps into a shared medium (CSMA/CD),
  collisions possible.
\item
  \textbf{Ring}: token-based, eliminates collisions, requires token
  management.
\item
  \textbf{Star}: master/slave, predictable, not scalable (master
  bottleneck).
\item
  \textbf{Mesh}: full or partial, robust, private, many physical links
\end{itemize}

\begin{center}\rule{0.5\linewidth}{0.5pt}\end{center}

\section{Endpoints}\label{endpoints}

\subsection{Sensors}\label{sensors}

\begin{itemize}
\tightlist
\item
  Devices detecting stimuli → electrical/processable output.
\item
  Types: pressure, temperature, biometric, chemical, acoustic, optical,
  positioning, etc.
\item
  Three classifications:

  \begin{itemize}
  \tightlist
  \item
    \textbf{Fixed}: stable, not moving like domotics
  \item
    \textbf{Mobile}: always moving (GPS, wearables)
  \item
    \textbf{Nomadic}: they move but then they stop somewhere else
    (portable medical devices, field sensors)
  \end{itemize}
\item
  \textbf{Active vs Passive} sensors

  \begin{itemize}
  \tightlist
  \item
    \textbf{Active}: need to run a current or send a signal to get a
    measurement (e.g., clocks, strain gauges, radar).
  \item
    \textbf{Passive}: use the physical phenomenon itself (e.g.,
    piezoaccelerometers, GPS, thermocouples).
  \end{itemize}
\end{itemize}

\subsubsection{Actuators}\label{actuators}

\begin{itemize}
\tightlist
\item
  Output elements performing actions.
\item
  \textbf{Digital}: on/off switching.
\item
  \textbf{Analog}: continuous signals (e.g., motor speed control).
\end{itemize}

\subsubsection{Gateways}\label{gateways}

\begin{itemize}
\tightlist
\item
  Bridge sensors ↔ higher-level processing/cloud.
\item
  Provide connectivity, data collection, filtering, storage, event
  processing, optional analytics
\end{itemize}

\subsubsection{Fog Nodes}\label{fog-nodes}

\begin{itemize}
\tightlist
\item
  ``Bring the cloud to the ground.''
\item
  More powerful than gateways, allow \textbf{local analytics with lower
  latency}.
\end{itemize}

\begin{center}\rule{0.5\linewidth}{0.5pt}\end{center}

\section{Data plane functions}\label{data-plane-functions}

\subsection{Data acquisition and
processing}\label{data-acquisition-and-processing}

\begin{itemize}
\tightlist
\item
  In order to get a digital signal from a sensor, we need to account
  for:

  \begin{itemize}
  \tightlist
  \item
    Sampling
  \item
    Aliasing
  \item
    Quantization
  \item
    Saturation
  \item
    Hysteresis and non-linearities
  \item
    Calibration
  \item
    Error propagation
  \end{itemize}
\end{itemize}

\begin{center}\rule{0.5\linewidth}{0.5pt}\end{center}

\subsubsection{Sampling}\label{sampling}

\begin{itemize}
\tightlist
\item
  First, we sample: we take (regular) measurements over time, and
  represent the signal with those measurements.
\item
  Nyquist theorem: if sampling is ≥ twice as fast as the signal,
  sampling has no errors.
\item
  Example: If a signal is thought to have a maximum frequency between
  1000 Hz and 4000 Hz, which of the following would be the most
  appropriate sample rate?

  \begin{itemize}
  \tightlist
  \item
    a. 500 Hz
  \item
    b. 8000 Hz
  \item
    c. 9000 Hz
  \item
    d. 24000 Hz
  \end{itemize}
\end{itemize}

\subsubsection{Aliasing}\label{aliasing}

\begin{itemize}
\tightlist
\item
  If we do not sample ``enough,'' we have aliasing: patterns in the
  sampled version.
\item
  A high frequency component in the frequency spectrum of a signal takes
  the identity of a lower frequency component in the same spectrum of
  the sampled signal.
\item
  This falsely estimated signal will be indistinguishable from (i.e., be
  an "alias" to) another signal having that true lower frequency.
\item
  There are tools to avoid this, but we need to know it is an issue.
\item
  Example: aliasing in a sine wave.\\
  \pandocbounded{\includegraphics[keepaspectratio]{/home/alessandrovulcu/Documenti/Università/ICT/Pasted image 20251007085600.png}}
\end{itemize}

\subsubsection{Quantization}\label{quantization}

\begin{itemize}
\tightlist
\item
  Digital values are finite: we cannot represent continuous values
  precisely!
\item
  Quantization always introduces an error, which is the sensor's natural
  error.
\item
  More bits, better resolution (but more data)
\end{itemize}

\pandocbounded{\includegraphics[keepaspectratio]{/home/alessandrovulcu/Documenti/Università/ICT/Pasted image 20251007085902.png}}

\subsubsection{Saturation}\label{saturation}

\begin{itemize}
\tightlist
\item
  Sensors often have an upper limit, which is either physical or
  dictated by the maximum quantization level: in this case, we have
  saturation (all higher values are mapped to the same output)
\item
  This is what happens when you take a picture with high exposure: all
  lighter points are bright white!\\
  \pandocbounded{\includegraphics[keepaspectratio]{/home/alessandrovulcu/Documenti/Università/ICT/Pasted image 20251007090152.png}}
\end{itemize}

\subsubsection{Hysteresis}\label{hysteresis}

\begin{itemize}
\item
  Sometimes sensors takes a while to track the process because the
  \textbf{physical nature of the measurement is slow}
\item
  Which means we have Hysteresis error when the delay of sensor
  measurements depends on the magnitude of the change.
\item
  Example: When you apply pressure to a sensor and then release that
  pressure, the sensor's output should in theory return you to the same
  value it started at before that specific pressure cycle. However, in
  reality, there are often small differences between the initial value
  and the final value after a pressure cycle.

  \begin{itemize}
  \tightlist
  \item
    From 0 to 100 bar
  \item
    0 bar → 0 V
  \item
    100 bar → 2500 V
  \item
    0 bar → 0.05 V

    \begin{itemize}
    \tightlist
    \item
      \textbf{0.02\% / 1 bar hysteresis}
    \end{itemize}
  \end{itemize}
\end{itemize}

\subsubsection{Non-linearities}\label{non-linearities}

\begin{itemize}
\item
  Ideally, sensors should also have linear responses: in reality, there
  are often non-linear effects that we need to adjust for.
\item
  Hard to predict
\item
  Example: A linear sensor's output will change by equal amounts for
  equal pressure changes across its entire range.

  \begin{itemize}
  \tightlist
  \item
    From 0 to 10 bar
  \item
    0 bar → 0 V
  \item
    10 bar → 5 V
  \item
    3 bar → 4.2 V (while I would have expected 3.75 V, according to the
    previous data).

    \begin{itemize}
    \tightlist
    \item
      \textbf{≤2\% non-linearity of full scale}
    \end{itemize}
  \end{itemize}
\end{itemize}

\subsubsection{Calibration}\label{calibration}

\begin{itemize}
\item
  Sensors often require calibration to give proper results:

  \begin{itemize}
  \tightlist
  \item
    Bias is a constant error term in one direction.
  \item
    Distortion is due to non-linearity.
  \end{itemize}
\item
  By measuring a series of known quantities, we can compensate for these
  issues (however, external factors such as temperature can throw off
  calibration!)
\item
  Example: Using the data from the calibration sheet, we see that the
  deviations are all less than the maximum deviation allowed of 0.5\%.
\end{itemize}

\subsubsection{Error propagation}\label{error-propagation}

\begin{itemize}
\tightlist
\item
  Calculations can change the error of a measurement, as the errors are
  also considered in the function you compute.
\item
  Non-linear functions can have weird effects with measurement errors.
\item
  Derivation (e.g., computing velocity from position measurements) is
  brittle: since the measurements are close together, noise can have a
  huge effect.
\item
  Integration/averaging are robust (rely on many points, and noise is
  compensated).
\item
  Example: Calculating heat index

  \begin{itemize}
  \tightlist
  \item
    Temperature sensor: measured (30°) with uncertainty ±0.5
  \item
    Humidity sensor: measured (70\%) with uncertainty ±2\% (= 0.02)
  \item
    Heat index = T + 0.33×RH−0.7 = 52.4
  \item
    \textbf{Uncertainty with error propagation: 0.828 \textgreater{}
    (0.5 + 0.02)}
  \end{itemize}
\end{itemize}

\subsubsection{External data
communication}\label{external-data-communication}

\begin{itemize}
\tightlist
\item
  Sensors data may need to be shared with external parties, such as
  other peer nodes at the edge, fog, and cloud levels.
\item
  Messages exchanged between IoT nodes may be delivered using
  \textbf{push} or \textbf{pull} modes.

  \begin{itemize}
  \tightlist
  \item
    \textbf{Pull}: a message is explicitly requested from the source
    (server) by its recipient, an IoT client.
  \item
    \textbf{Push}: the data source and the gateway send data when the
    specified conditions are met, such as arrival of the data sampling
    time mark, or exceeding of the measurement threshold.
  \end{itemize}
\item
  \textbf{Publish/subscribe system}: data sources ``publish'' data to a
  known place, and authorized clients can ``subscribe'' to data of
  interest and receive related messages.

  \begin{itemize}
  \tightlist
  \item
    Publishers do not need to know the identity or number of their
    receivers (subscribers), and the receivers do not need to know the
    identity or state of their publishers.
  \end{itemize}
\end{itemize}

\paragraph{How to share sensors' data?}\label{how-to-share-sensors-data}

\begin{itemize}
\tightlist
\item
  Some processes (e.g., an electricity meter, a room thermostat) do not
  need to be sampled very often.
\item
  If we only care about cumulative measures (how much electricity was
  consumed in total), we do not need frequent samples!
\item
  On the other hand, faster processes need frequent samples, as process
  can change faster.
\item
  Why don't we transmit all the time?
\item
  There are several reasons why we cannot transmit all the time, but we
  will concentrate on two of them in this first part of the course:

  \begin{itemize}
  \tightlist
  \item
    Limited bandwidth
  \item
    Limited energy
  \end{itemize}
\item
  Each transmission (and, in fact, each activity the sensor performs)
  costs energy, as well as blocking the wireless channel for other
  transmissions
\end{itemize}

\paragraph{Some solutions to send sensors' data more
efficiently}\label{some-solutions-to-send-sensors-data-more-efficiently}

\begin{itemize}
\item
  \textbf{Interpolation}

  \begin{itemize}
  \tightlist
  \item
    Interpolation is the mathematical tool to infer missing data from
    sparse samples. We have to make some assumptions on the nature of
    the process to be able to interpolate: if we have the correct model,
    we can send a lot fewer data samples!
  \item
    Interpolation can also be dangerous: we can underfit (select a model
    that is too simple) or overfit (select a model that is too complex
    and ends up following random noise).\\
    \pandocbounded{\includegraphics[keepaspectratio]{/home/alessandrovulcu/Documenti/Università/ICT/Pasted image 20251007093806.png}}
  \end{itemize}
\item
  \textbf{Spatial correlation}

  \begin{itemize}
  \tightlist
  \item
    Data are usually correlated in space:

    \begin{itemize}
    \tightlist
    \item
      Traffic on either side of a bridge.
    \item
      Temperatures close together.
    \end{itemize}
  \item
    There can be easy natural relations (points close together are
    similar) or structural relations (connected points are similar).
  \item
    Since measurements have errors, correlations help us improving the
    quality of estimates: we can make statistical calculations and
    estimates to figure out what is happening in points (like the purple
    one on the bridge) where we have no sensors, or refine our
    estimates.
  \end{itemize}
\item
  \textbf{Temporal correlation}

  \begin{itemize}
  \tightlist
  \item
    Data are also correlated in time:

    \begin{itemize}
    \tightlist
    \item
      Traffic in one spot
    \item
      Temperature measured by one sensor
    \end{itemize}
  \item
    The relation here can be more or less complex, and depend on space
    as well (e.g., line at a traffic light), but we can also exploit
    this to get better data with fewer transmissions.
  \item
    Correlations can be misleading and dangerous! When relying on
    statistical information to reconstruct missing data or make
    decisions, we need to be aware that random chance is always a
    possibility.
  \end{itemize}
\end{itemize}

\begin{center}\rule{0.5\linewidth}{0.5pt}\end{center}

\subsubsection{Local control}\label{local-control}

\begin{itemize}
\tightlist
\item
  Provided in systems where latency is critical and/or edge autonomy is
  required.
\item
  Local control enables the edge node to execute pre-defined control
  sequences when a specified condition occurs on some combination of
  local data.
\item
  If This Then That: scenes where lights and heating/cooling are
  adjusted accordingly when users are leaving or approaching their
  homes.
\item
  \textbf{Node-RED}: provides a graphical user interface to construct
  data flows with the functional processing modules and actuation output
  stages.
\item
  In the extreme case of temperature, the script triggers the device
  shutdown.
\end{itemize}

\pandocbounded{\includegraphics[keepaspectratio]{/home/alessandrovulcu/Documenti/Università/ICT/Pasted image 20251007094532.png}}

\subsubsection{Data storage}\label{data-storage}

\begin{itemize}
\tightlist
\item
  Data storage for the acquired data may be provided by the gateway.
\item
  The existence of local storage provides the flexibility to conserve
  network bandwidth and to reduce the load on cloud resources by
  selectively reporting data that meet certain conditions.

  \begin{itemize}
  \tightlist
  \item
    \textbf{Semantic communication.}
  \item
    Depending on the nature of the physical thing being measured,
    \textbf{only some of the data points may be of interest}. E.g.,
    successive minor changes may be of little interest, but those that
    exceed a comfort guard band or change by a significant percentage or
    amount from prior readings are.
  \item
    Qualifying readings may be defined as events and reported only when
    they occur, with raw readings stored in the local database should
    they be required for auditing or subsequent analysis.
  \end{itemize}
\end{itemize}

\subsubsection{(Edge) analytics}\label{edge-analytics}

\begin{itemize}
\tightlist
\item
  Edge analytics for the data may be provided by the gateway or the fog
  node.
\item
  Proximity to data sources: low latency.
\item
  Ability to process high-frequency data samples (at the rate of the
  sensor) → real-time.
\item
  Conservation of network bandwidth by not sending data to the cloud.
\item
  Can operate even in disconnected mode.
\item
  Analytics generally require AI/ML models:

  \begin{itemize}
  \tightlist
  \item
    Insufficient computational resources at the edge.
  \item
    No data aggregation of edge/local data.
  \end{itemize}
\item
  IoT systems should be implemented using design tools and practices
  that facilitate flexible allocation of functions (in tandem between
  edge and cloud).
\end{itemize}

\subsubsection{Analytics placement}\label{analytics-placement}

\begin{itemize}
\tightlist
\item
  Optimal placement of data and processing functions is an age-old
  trade-off in distributed systems.
\item
  It is based on whether it is more cost-effective to move the data to
  the computation or to move the computation to the data, depending on:

  \begin{itemize}
  \tightlist
  \item
    Availability and cost of bandwidth
  \item
    Latency requirements for time-critical operations
  \item
    Local autonomy of operations, including disconnected mode
  \item
    Security and data control or privacy concerns
  \item
    Cost and complexity of managing distributed nodes with computing and
    storage
  \item
    Available energy resources
  \item
    Available computational capacity
  \item
    \ldots{}
  \end{itemize}
\end{itemize}

\begin{center}\rule{0.5\linewidth}{0.5pt}\end{center}

\section{Control plane functions}\label{control-plane-functions}

\begin{itemize}
\tightlist
\item
  Fully functional gateways may be fairly sophisticated computational
  and communications nodes that host and execute a variety of services.
\item
  They need to be properly secured and managed.
\item
  Main control functionalities:

  \begin{itemize}
  \tightlist
  \item
    Fault management (troubleshooting, error logging, and recovery).
  \item
    Remote monitoring, control, administration, and diagnostics.
  \item
    Remote firmware and software updates.
  \item
    Security updates.
  \item
    Metering (network bandwidth and software usage) and supervision.
  \item
    Provisioning and authentication.
  \end{itemize}
\end{itemize}

\begin{center}\rule{0.5\linewidth}{0.5pt}\end{center}

\subsubsection{(Cyber)security}\label{cybersecurity}

\begin{itemize}
\item
  The first step is to determine the security risk.
\item
  Identify, quantify, and prioritize the risks for the system, and find
  a balance between costs and potential losses from residual risks.
\item
  Model to quantify risk: vulnerability, threats, and impacts.\\
  \pandocbounded{\includegraphics[keepaspectratio]{/home/alessandrovulcu/Documenti/Università/ICT/Pasted image 20251007095925.png}}
\item
  \textbf{Vulnerabilities}: can originate from the design of the system
  architecture, technology, and protocols, but more often emerge in the
  implementation (e.g., coding flaws) and in the operation phases (e.g.,
  configuration errors, use of weak credential, unparched devices,
  unexpert use, etc.).\\
\item
  \textbf{Security measures for the diligent user}:

  \begin{itemize}
  \tightlist
  \item
    Partitioning the system into zones (characterized in terms of
    security levels).
  \item
    Adopt effective authentication protocols, tailored to the different
    zones.
  \item
    Provide cryptographic protection of data.
  \item
    Install equipment from trusted providers.
  \item
    Make and install frequent patches and workarounds (software
    updates).
  \item
    \ldots{}
  \end{itemize}
\item
  \textbf{Threat}: an IoT network is a possible point of attack if the
  endpoint is unguarded and is positioned in an easily reachable
  position.

  \begin{itemize}
  \tightlist
  \item
    \textbf{Eavesdropping}: learning confidential information.
  \item
    \textbf{Traffic analysis}: learning data origin, destination, and
    length (not the content).
  \item
    \textbf{Forgery}: building a fake message, pretending it was sent by
    someone else.
  \item
    \textbf{Masquerade}: claiming to be someone else in a single
    message.
  \item
    \textbf{Repudiation}: denying having sent or received a message.
  \item
    \textbf{Profiling}: gathering information about a single user.
  \item
    \textbf{Fingerprinting}: identifying the user associated with a
    message.
  \item
    \textbf{Jamming}: causing interference to a communication.
  \end{itemize}
\item
  \textbf{Standards}: IoT-oriented and industry-specific standards have
  been developed for security.\\
  \pandocbounded{\includegraphics[keepaspectratio]{/home/alessandrovulcu/Documenti/Università/ICT/Pasted image 20251007100557.png}}
\end{itemize}

\begin{center}\rule{0.5\linewidth}{0.5pt}\end{center}

\subsubsection{Provisioning and
authentication}\label{provisioning-and-authentication}

\begin{itemize}
\tightlist
\item
  \textbf{Provisioning}: the lifecycle of an endpoint begins with its
  provisioning in the network.
\item
  It is the process of giving the endpoint the configuration data and
  the ability to present itself.
\item
  \textbf{Authorization}: checking the rights of access to the network.
\item
  \textbf{Application-level authentication}: used if the network is
  shared among applications.
\item
  \textbf{Authentication}: verification of the identity (``I know your
  name, is it you?'').
\item
  \textbf{Challenge-response protocol}: employs secret keys for
  verification.
\end{itemize}

\pandocbounded{\includegraphics[keepaspectratio]{/home/alessandrovulcu/Documenti/Università/ICT/Pasted image 20251007101013.png}}

\begin{center}\rule{0.5\linewidth}{0.5pt}\end{center}

\section{QoS and performance}\label{qos-and-performance}

\begin{itemize}
\tightlist
\item
  \textbf{One-way delay (OWD)}: time between the transmission of the
  first bit of a packet at an interface and the reception of the last
  bit of the packet at the other interface.
\item
  Requires perfect time synchronization between the endpoints.
\item
  \textbf{Time reference}: Network Time Protocol (NTP).
\item
  \textbf{Two-way delay (TWD)}: time between the transmission of the
  first bit of a packet at an interface and the reception of the last
  bit of the packet at the same interface, after being reflected by the
  system at the other end, net of the End System Delay (ESD, processing
  time at the remote node).
\item
  We have that \textbf{Round Trip Time (RTT) = TWD + ESD.}
\end{itemize}

\subsection{Delay}\label{delay}

\begin{itemize}
\item
  \textbf{Packet-delay variation (PDV)} (\textasciitilde jitter): based
  on previous OWD/TWD measures.

  \begin{itemize}
  \tightlist
  \item
    PDV = \textbar OWDi - OWDi-1\textbar{}
  \item
    PDV = \textbar OWDi - mean\{OWD\}\textbar{}
  \item
    PDV = \textbar OWDi - min\{OWD\}\textbar{}
  \item
    PDV = \textbar OWDi - OWD1\textbar{}
  \item
    \ldots{}
  \end{itemize}
\item
  \textbf{Problem}: These delay metrics do not consider the MAC delay
  (for channel access).
\item
  \textbf{End-to-end delay (E2E)}: it consists of (at least) three
  components:

  \begin{itemize}
  \tightlist
  \item
    \textbf{MAC delay} (under the control of the system designer).
  \item
    \textbf{Access link and gateway delay} (under the control of the
    system designer).
  \item
    \textbf{OWD} (depends on the WAN).
  \end{itemize}
\end{itemize}

\subsubsection{Delay and ``real-time''}\label{delay-and-real-time}

\begin{itemize}
\tightlist
\item
  The expression \textbf{``real-time''} may have different
  interpretations:

  \begin{itemize}
  \tightlist
  \item
    \textbf{Synchronization} to a common time reference.
  \item
    \textbf{Responding fast.}
  \item
    \textbf{Responding within predictable response times.}\\
  \end{itemize}
\item
  \textbf{Hard real-time}: if the response exceeds a limit, the system
  does not operate well.

  \begin{itemize}
  \tightlist
  \item
    Example: respond within 5 ms
  \item
    \textbf{CASE 1:} The application always responds in 4 ms

    \begin{itemize}
    \tightlist
    \item
      Average delay: 4 ms; Failure rate: 0\%
    \end{itemize}
  \item
    \textbf{CASE 2:} The application responds in 1 ms for 99.9\% of the
    times, and in 6 ms in 0.01\%.

    \begin{itemize}
    \tightlist
    \item
      Average delay: 1.005 ms; Failure rate: 10⁻³
    \end{itemize}
  \end{itemize}
\item
  \textbf{Soft real-time}: it involves some response time percentiles.
\end{itemize}

\subsubsection{Packet loss}\label{packet-loss}

\begin{itemize}
\tightlist
\item
  In an IoT network, \textbf{loss may occur} if:

  \begin{itemize}
  \tightlist
  \item
    The \textbf{buffers of a node overflow.}
  \item
    \textbf{Transmission errors} (negligible in wired LANs, dominant in
    wireless LANs).
  \item
    \textbf{Delay much larger than expected} may be considered as a
    loss.\\
  \end{itemize}
\item
  How to have realistic measures of packet loss? There are some
  problems:

  \begin{itemize}
  \tightlist
  \item
    \textbf{Long-lasting measurements:} averaging of data gives no
    indications on the short period.
  \item
    \textbf{Short measurements:} no statistically meaningful data.
  \item
    \textbf{Intense bursts:} too intrusive, altering the result (not
    representative of a real situation).
  \end{itemize}
\end{itemize}

\subsubsection{Capacity}\label{capacity}

\begin{itemize}
\tightlist
\item
  \textbf{Hardly a problem}, as applications are not (in general)
  bandwidth demanding.
\item
  For some high-end applications (e.g., telemedicine, IIoT, automotive,
  etc.), \textbf{capacity} is generally less important than other
  metrics such as reliability.
\item
  \textbf{Example:} metering application

  \begin{itemize}
  \tightlist
  \item
    Collect 1M samples (or sensors?) (of 128B) every 15 minutes.
  \item
    \textbf{Average throughput:} 1,000,000 × 128B × 8 (bit) / (15 × 60)
    = 1.3 Mbit/s
  \item
    If the same measurements take place as unsolicited events from the
    endpoint in the same second, the \textbf{link capacity} shall be
    around 1,000,000 × 128B × 8 (bit) = 1 Gbit/s.
  \end{itemize}
\end{itemize}

Those are the most important metrics, but there are other ones related
to IoT.

\subsubsection{Reliability}\label{reliability}

\begin{itemize}
\tightlist
\item
  \textbf{R(t)}: probability that a system is working as expected at a
  certain time \emph{t.}

  \begin{itemize}
  \tightlist
  \item
    For most systems, the early phase of ``high mortality'' is short and
    occurs during tests.
  \item
    The system is generally replaced for obsolescence as the aging phase
    approaches.
  \item
    In ``normal life'', most systems shall be reliable.
  \end{itemize}
\item
  \textbf{R(t)} is memoryless: the residual life of failure does not
  depend on how long the system has been working → the system does not
  age during ``normal life.''\\
  \pandocbounded{\includegraphics[keepaspectratio]{/home/alessandrovulcu/Documenti/Università/ICT/Pasted image 20251008085321.png}}
\end{itemize}

\subsubsection{Availability}\label{availability}

\begin{itemize}
\tightlist
\item
  \textbf{Probability that the system is working} when observed at a
  random time instant.
\item
  We assume that the system can be periodically repaired if needed and
  put back in service.
\item
  \textbf{Mean Time Between Failures (MTBF).}
\item
  \textbf{Mean Time To Repair (MTTR).}
\end{itemize}

\pandocbounded{\includegraphics[keepaspectratio]{/home/alessandrovulcu/Documenti/Università/ICT/Pasted image 20251008085509.png}}

\paragraph{Availability (example)}\label{availability-example}

\begin{itemize}
\tightlist
\item
  For memoryless systems, \textbf{MTBF = Mean Time To Fail (MTTF) =
  1/λ.}\\
\item
  \textbf{Example:} prevention of fires in a forest

  \begin{itemize}
  \tightlist
  \item
    \emph{N} sensors that monitor fires in the forest.
  \item
    Sensors fail at rate \textbf{λ.}
  \item
    The prevention works if \emph{x\%} of sensors are active and alive.
  \item
    Period of effectiveness of the system (\textbf{T}):
  \end{itemize}
\end{itemize}

\begin{center}\rule{0.5\linewidth}{0.5pt}\end{center}

\section{Energy constraints}\label{energy-constraints}

\subsection{Introduction}\label{introduction}

\begin{itemize}
\tightlist
\item
  Most sensors and components use electrical energy, i.e., they exploit
  the movements of electrons across a conducting medium.
\item
  Powering is a very important issue, and depends on the endpoint
  capabilities and the network architecture procedures.
\end{itemize}

\begin{longtable}[]{@{}ll@{}}
\toprule\noalign{}
Type & Examples \\
\midrule\noalign{}
\endhead
\bottomrule\noalign{}
\endlastfoot
\textbf{Mains / manually recharged} & Field buses Smartphones Smart grid
devices \\
\textbf{Battery} & Most sensors \\
\textbf{Data line} & Videocameras \\
\textbf{Harvesting} & Environmental sensors \\
\end{longtable}

\begin{center}\rule{0.5\linewidth}{0.5pt}\end{center}

\subsection{Sources of consumption}\label{sources-of-consumption}

\subsubsection{Computing}\label{computing}

\begin{itemize}
\tightlist
\item
  Computing requires a certain energy for each cycle, corresponding to
  one basic operation. Some operations (e.g., sinusoidal functions)
  require multiple cycles.

  \begin{itemize}
  \tightlist
  \item
    Energy consumption is usually a function of the number of cycles
    required.
  \end{itemize}
\item
  If a device has a certain clock frequency, we can also compute the
  power it requires to perform calculations.
\item
  Computing energy can be formalized based on the number of logical
  operations that are required to perform: \textbf{AND, NOT, OR, XOR.}
\end{itemize}

\begin{center}\rule{0.5\linewidth}{0.5pt}\end{center}

\subsubsection{Boards and PLCs}\label{boards-and-plcs}

\begin{itemize}
\tightlist
\item
  Boards and PLCs are standard circuits used in industrial automation.
  Larger IoT nodes also use Arduinos or other pre-made boards.
\item
  Smaller, more efficient nodes are often designed ad hoc, to use as
  little energy as possible. In general, IoT nodes have limited
  computational capabilities.
\end{itemize}

\textbf{Example: Arduino-UNO}\\
The Arduino Uno is a microcontroller board based on the ATmega328.\\
It contains everything needed to support the microcontroller; simply
connect it to a computer with a USB cable or power it with an AC-to-DC
adapter or battery to get started.

\begin{center}\rule{0.5\linewidth}{0.5pt}\end{center}

\subsubsection{Communication}\label{communication}

\begin{itemize}
\tightlist
\item
  Digital communication usually requires more steps than just sending a
  wave, and each consumes energy:
\end{itemize}

\begin{enumerate}
\tightlist
\item
  \textbf{Encoding and error protection.}

  \begin{itemize}
  \tightlist
  \item
    Use known redundant codes: the simplest example is the parity check
    bit.
  \item
    Decoders are often complex, but can correct errors in transmission
    and recover the original packet. However, complexity is paid for in
    energy.\\
  \end{itemize}
\item
  \textbf{Preambles for synchronization and channel estimation.}
\item
  \textbf{Modulation} from the baseband digital signal to the passband
  analog signal.
\item
  \textbf{Transmission} (power is applied to the antenna) (not required
  for reception).

  \begin{itemize}
  \tightlist
  \item
    IoT transmissions are usually slow (low bitrate), but distance plays
    a large role.
  \end{itemize}
\end{enumerate}

\begin{center}\rule{0.5\linewidth}{0.5pt}\end{center}

\subsection{Source of energy}\label{source-of-energy}

\subsubsection{Mains power}\label{mains-power}

\begin{itemize}
\tightlist
\item
  Provide virtually unlimited power, making it possible to employ
  complex protocols, work at high bitrates, and achieve long-distance
  communication.
\item
  There are some drawbacks:

  \begin{itemize}
  \tightlist
  \item
    Only work in wired fixed systems.
  \item
    Can be dangerous in environments at risk of fire and explosions.
  \item
    Need to lay out an electric line (challenging in hard-to-reach
    remote/rural areas).
  \item
    Need to put transformers (expensive).
  \end{itemize}
\end{itemize}

\begin{center}\rule{0.5\linewidth}{0.5pt}\end{center}

\subsubsection{Power over Ethernet (PoE)}\label{power-over-ethernet-poe}

\begin{itemize}
\tightlist
\item
  Employ the wired communication infrastructure also for powering.
\item
  Can carry DC power on the twisted pair used for Ethernet connectivity.
\item
  Only one cabled infrastructure needs to be in place.
\item
  Transformers are not needed.
\end{itemize}

\begin{center}\rule{0.5\linewidth}{0.5pt}\end{center}

\subsubsection{Powerline communication}\label{powerline-communication}

\begin{itemize}
\tightlist
\item
  It is the opposite of PoE.
\item
  Use the powerline to carry data.
\item
  Based on modulation between high-frequency data and low-frequency
  electricity.
\item
  Only one cabled infrastructure needs to be in place.
\item
  Advantageous when the system to be monitored and controlled is the
  power grid itself.\\
  \pandocbounded{\includegraphics[keepaspectratio]{/home/alessandrovulcu/Documenti/Università/ICT/Pasted image 20251008091052.png}}
\end{itemize}

\begin{center}\rule{0.5\linewidth}{0.5pt}\end{center}

\subsubsection{Batteries}\label{batteries}

\begin{itemize}
\tightlist
\item
  Batteries use the charge difference between ions to induce a potential
  difference.

  \begin{itemize}
  \tightlist
  \item
    \textbf{Rechargeable} batteries can run the circuit backwards and
    move the electrons back (by connecting a stronger generator).
  \item
    Charging is never perfect: chemical impurities build up.
  \end{itemize}
\item
  If the battery runs out completely, it is often impossible to recharge
  without direct intervention (which may be expensive): the objective is
  \textbf{never to let it run out!}

  \begin{itemize}
  \tightlist
  \item
    Apply \textbf{statistical data} and direct information from the
    meter to define a plan of \textbf{proactive} on-field
    \textbf{intervention} if needed → carry out with proper timing to
    minimize the number of substitutions in the lifetime and the number
    of dedicated ``truck rolls'' or manual readings.
  \end{itemize}
\end{itemize}

\subsubsection{Batteries (management and
optimization)}\label{batteries-management-and-optimization}

\begin{itemize}
\tightlist
\item
  Energy is the basic constraint of IoT systems. What can we do?

  \begin{itemize}
  \tightlist
  \item
    Employ battery technology with low self-discharging properties.
  \item
    Minimize the duration of active states of the radio components with
    suitable scheduling algorithms and MAC protocols.

    \begin{itemize}
    \tightlist
    \item
      Duty cycles, sleep mode, alert/notification, etc.
    \end{itemize}
  \item
    Adopt energy-aware routing.
  \item
    Introduce endpoint management protocols that communicate the state
    of batteries.
  \item
    Adopt a ``let it be off'' approach.
  \item
    Recharge batteries with energy harvesting.
  \end{itemize}
\end{itemize}

\begin{center}\rule{0.5\linewidth}{0.5pt}\end{center}

\subsubsection{Energy harvesting}\label{energy-harvesting}

\begin{itemize}
\tightlist
\item
  Energy harvesting: recharge it while it is running.
\item
  Energy harvesting sensors and actuators are equipped with some kind of
  generator and can recharge their battery from it. This allows them to
  stay alive much longer than with a single battery charge, even if they
  get very little energy.\\
  \pandocbounded{\includegraphics[keepaspectratio]{/home/alessandrovulcu/Documenti/Università/ICT/Pasted image 20251008091530.png}}
\end{itemize}

\begin{center}\rule{0.5\linewidth}{0.5pt}\end{center}

\subsubsection{Solar and wind energy}\label{solar-and-wind-energy}

\begin{itemize}
\tightlist
\item
  \textbf{Solar panels} are now cheap, and solar energy is renewable and
  clean.

  \begin{itemize}
  \tightlist
  \item
    The sensor needs to be exposed (shade can be a problem, underground
    or indoor sensors will not work).
  \item
    Photovoltaic cells will not generate any power at night (the sensor
    needs to get through the night on a single charge).
  \end{itemize}
\item
  Another possibility is wind energy, but there are several challenges:

  \begin{itemize}
  \tightlist
  \item
    The sensor still needs to be exposed (underground or indoor sensors
    will not work).
  \item
    Wind is intermittent and unpredictable.
  \item
    Small turbines are inefficient.
  \end{itemize}
\end{itemize}

\begin{center}\rule{0.5\linewidth}{0.5pt}\end{center}

\subsubsection{Thermal and vibration
energy}\label{thermal-and-vibration-energy}

\begin{itemize}
\tightlist
\item
  \textbf{Thermal energy} is cheap when the sensor has access to
  temperature gradients:

  \begin{itemize}
  \tightlist
  \item
    The inside/outside or ground/air gradient can generate energy.
  \item
    The power is usually very small, but may be the only option for
    buried sensors.
  \item
    Need for long distances between terminals.
  \item
    Good for wearables: the difference between body temperature and the
    air is high.
  \end{itemize}
\item
  Vibration energy can be collected using piezoelectric devices with
  good efficiency.
\item
  The constant presence of vibrations energetic enough to power the
  endpoint is uncertain, except in some industrial environments.
\end{itemize}

\begin{center}\rule{0.5\linewidth}{0.5pt}\end{center}

\subsubsection{Wireless power transfer
(WPT)}\label{wireless-power-transfer-wpt}

\begin{itemize}
\tightlist
\item
  Sensors can benefit from the power of received transmissions, but the
  benefits are tiny at longer distances: this is only feasible for very
  low-power sensors.
\item
  \textbf{Inductive coupling}: uses electromagnetic fields to transfer
  energy between coils.
\item
  \textbf{Capacitive coupling}: transfers energy via electric fields
  between conductive plates.
\item
  Uses \textbf{RF signals} to transfer power over longer distances
  (several meters to kilometers).
\item
  Uses \textbf{focused laser beams} to transmit power over long
  distances.
\end{itemize}

\begin{center}\rule{0.5\linewidth}{0.5pt}\end{center}

\subsubsection{Passive sensors}\label{passive-sensors}

\begin{itemize}
\tightlist
\item
  We can also use the received transmission to power the circuit
  directly: \textbf{RFID tags, biometric passports, and contactless
  cards} work this way.
\item
  Power loss is huge: these only work over short distances.
\item
  The sensor cannot actively transmit or record data, only respond to
  messages.
\item
  There are \textbf{privacy issues} (anyone can read a tag, the only
  limit is the SNR).\\
  \pandocbounded{\includegraphics[keepaspectratio]{/home/alessandrovulcu/Documenti/Università/ICT/Pasted image 20251008091951.png}}
\end{itemize}

\end{document}
